% Part: first-order-logic
% Chapter: natural-deduction
% Section: proving-things

\documentclass[../../../include/open-logic-section]{subfiles}

\begin{document}

\iftag{FOL}
      {\olfileid{fol}{ntd}{pro}}
      {\olfileid{pl}{ntd}{pro}}

\olsection{Voorbeelden van \usetoken{P}{derivation}}

\begin{ex}
Laten we !!a{derivation} geven van de !!{sentence}
$(!A \land !B) \lif !A$.

We beginnen door de gewenste conclusie onderaan de !!{derivation} te
schrijven.
\begin{prooftree}
\AxiomC{}
\UnaryInfC{$(!A\land !B) \lif !A$}
\end{prooftree}

Vervolgens moeten we bepalen welk soort inferentie !!a{sentence} van
deze vorm als resultaat kan hebben. De !!{main operator} van de
conclusie is $\lif$. We proberen de conclusie daarom met de regel
\Intro{\lif} te verkrijgen. Het is raadzaam de betrokken aannamen op
te schrijven en de afleidingsregels gaandeweg te benoemen. Zo is
gemakkelijk te zien of aan het einde van het bewijs alle aannamen zijn
!!{discharged}.
\begin{prooftree}
\AxiomC{$\Discharge{!A \land !B}{1}$}
\DeduceC{$!A$}
\DischargeRule{\Intro{\lif}}{1} 
\UnaryInfC{$(!A\land !B) \lif !A$}
\end{prooftree}

Nu moeten we de stappen van de aanname $!A \land !B$ naar $!A$
invullen. Omdat we slechts met één connectief te maken hebben,
$\land$, moeten we de $\land$-eliminatieregel gebruiken. Dat levert
het volgende bewijs op:
\begin{prooftree}
\AxiomC{$\Discharge{!A \land !B}{1}$}
\RightLabel{\Elim{\land}}
\UnaryInfC{$!A$}
\DischargeRule{\Intro{\lif}}{1} 
\UnaryInfC{$(!A\land !B) \lif !A$}
\end{prooftree}
We hebben nu een correcte !!{derivation} van $(!A \land
!B) \lif !A$.
\end{ex}

\begin{ex}
Laten we nu !!a{derivation} geven van $(\lnot !A \lor !B)
\lif (!A \lif !B)$.

We beginnen door de gewenste conclusie onderaan de !!{derivation} te
schrijven.
\begin{prooftree}
\AxiomC{}
\UnaryInfC{$(\lnot !A \lor !B) \lif (!A \lif !B)$}
\end{prooftree}
Om een logische regel te vinden die deze conclusie kan opleveren,
bekijken we de logische connectieven in de conclusie: $\lnot$,
$\lor$ en $\lif$. Voorlopig is alleen het eerste voorkomen van
$\lif$ van belang, want dat is de !!{main operator} van de
!!{sentence} in het eindsequent. $\lnot$, $\lor$ en het tweede
voorkomen van $\lif$ vallen binnen het bereik van een ander
connectief; die behandelen we dus later. We beginnen daarom met de
regel \Intro{\lif}. Een correcte toepassing moet er als volgt uitzien:
\begin{prooftree}
\AxiomC{$\Discharge{\lnot !A \lor !B}{1}$}
\DeduceC{$!A \lif !B$}
\DischargeRule{\Intro{\lif}}{1}
\UnaryInfC{$(\lnot !A \lor !B) \lif (!A \lif !B)$}
\end{prooftree}
Er zijn nu twee manieren om verder te gaan. We kunnen van onder naar
boven blijven werken en naar een tweede toepassing van de regel
\Intro{\lif} zoeken, of we kunnen van boven naar beneden werken en een
regel \Elim{\lor} toepassen. Laten we voor het laatste kiezen. We
gebruiken de aanname $\lnot !A \lor !B$ als de meest linkse premisse
van \Elim{\lor}. Voor een geldige toepassing van \Elim{\lor} moeten
de twee andere premissen identiek zijn aan de conclusie
$!A \lif !B$. Elk van die premissen mag achtereenvolgens uit een
andere aanname worden afgeleid, namelijk uit een van de twee disjuncten
van $\lnot !A \lor !B$. Onze !!{derivation} ziet er dus als volgt uit:
\begin{prooftree}
\AxiomC{$\Discharge{\lnot !A \lor !B}{1}$}
\AxiomC{$\Discharge{\lnot !A}{2}$}
\DeduceC{$!A \lif !B$}
\AxiomC{$\Discharge{!B}{2}$}
\DeduceC{$!A \lif !B$}
\DischargeRule{\Elim{\lor}}{2}
\TrinaryInfC{$!A \lif !B$}
\DischargeRule{\Intro{\lif}}{1} 
\UnaryInfC{$(\lnot !A \lor !B) \lif (!A \lif !B)$}
\end{prooftree}

In elk van de twee takken rechts willen we $!A \lif !B$ !!{derive}.
Dat kan het best met \Intro{\lif}.
\begin{prooftree}
\AxiomC{$\Discharge{\lnot !A \lor !B}{1}$}
\AxiomC{$\Discharge{\lnot !A}{2}, \Discharge{!A}{3}$}
\DeduceC{$!B$}
\DischargeRule{\Intro{\lif}}{3}
\UnaryInfC{$!A \lif !B$}
\AxiomC{$\Discharge{!B}{2}, \Discharge{!A}{4}$}
\DeduceC{$!B$}
\DischargeRule{\Intro{\lif}}{4}
\UnaryInfC{$!A \lif !B$}
\DischargeRule{\Elim{\lor}}{2}
\TrinaryInfC{$!A \lif !B$}
\DischargeRule{\Intro{\lif}}{1} 
\UnaryInfC{$(\lnot !A \lor !B) \lif (!A \lif !B)$}
\end{prooftree}

Voor de twee ontbrekende delen van de !!{derivation} hebben we in het
midden !!{derivation}s van $!B$ uit $\lnot !A$ en $!A$ nodig, en
links uit $!A$ en $!B$. Laten we het eerste geval eerst
behandelen. $\lnot !A$ en $!A$ zijn de twee premissen van
\Elim{\lnot}:
\begin{prooftree}
\AxiomC{$\Discharge{\lnot !A}{2}$}
\AxiomC{$\Discharge{!A}{3}$}
\RightLabel{\Elim{\lnot}}
\BinaryInfC{$\lfalse$}
\DeduceC{$!B$}
\end{prooftree}
Met \FalseInt kunnen we $!B$ als conclusie verkrijgen en de tak
voltooien.
\begin{prooftree}
\AxiomC{$\Discharge{\lnot !A \lor !B}{1}$}
\AxiomC{$\Discharge{\lnot !A}{2}$}
\AxiomC{$\Discharge{!A}{3}$}
\RightLabel{\Intro{\lfalse}}
\BinaryInfC{$\lfalse$}
\RightLabel{\FalseInt}
\UnaryInfC{$!B$}
\DischargeRule{\Intro{\lif}}{3}
\UnaryInfC{$!A \lif !B$}
\AxiomC{$\Discharge{!B}{2}, \Discharge{!A}{4}$}
\DeduceC{$!B$}
\DischargeRule{\Intro{\lif}}{4}
\UnaryInfC{$!A \lif !B$}
\DischargeRule{\Elim{\lor}}{2}
\TrinaryInfC{$!A \lif !B$}
\DischargeRule{\Intro{\lif}}{1} 
\UnaryInfC{$(\lnot !A \lor !B) \lif (!A \lif !B)$}
\end{prooftree}

Laten we nu de meest rechtse tak bekijken. Hier is het van belang te
beseffen dat de definitie van !!{derivation} \emph{toestaat dat
aannamen worden opgeheven}, maar dit \emph{niet vereist}. Als we
$!B$ uit een van de aannamen $!A$ en $!B$ kunnen afleiden zonder de
andere te gebruiken, is dat dus toegestaan. $!B$ uit~$!B$ !!{derive}
is triviaal: $!B$ op zichzelf is al zo'n !!{derivation}, zonder dat
inferenties nodig zijn. We kunnen de aanname~$!A$ daarom eenvoudig
weglaten.
\begin{prooftree}
\AxiomC{$\Discharge{\lnot !A \lor !B}{1}$}
\AxiomC{$\Discharge{\lnot !A}{2}$}
\AxiomC{$\Discharge{!A}{3}$}
\RightLabel{\Elim{\lnot}}
\BinaryInfC{$\lfalse$}
\RightLabel{\FalseInt}
\UnaryInfC{$!B$}
\DischargeRule{\Intro{\lif}}{3}
\UnaryInfC{$!A \lif !B$}
\AxiomC{$\Discharge{!B}{2}$}
\RightLabel{\Intro{\lif}}
\UnaryInfC{$!A \lif !B$}
\DischargeRule{\Elim{\lor}}{2}
\TrinaryInfC{$!A \lif !B$}
\DischargeRule{\Intro{\lif}}{1} 
\UnaryInfC{$(\lnot !A \lor !B) \lif (!A \lif !B)$}
\end{prooftree}
Merk op dat de meest rechtse \Intro{\lif}-inferentie in de voltooide
!!{derivation} feitelijk geen enkele aanname opheft.
\end{ex}

\begin{ex}
Tot dusver hebben we de regel \FalseCl{} niet nodig gehad. Deze regel
is bijzonder omdat we er een aanname mee mogen opheffen die geen
sub-!!{formula} van de conclusie van de regel is. Hij hangt nauw samen
met de regel \FalseInt{}. De regel \FalseInt{} is zelfs een bijzonder
geval van de regel \FalseCl{}. Er bestaat een logica die
‘intuïtionistische logica’ heet en waarin alleen \FalseInt{} is
toegestaan. De regel \FalseCl{} is een laatste
redmiddel wanneer niets anders werkt. Stel bijvoorbeeld dat we
$!A \lor \lnot !A$ willen !!{derive}. Onze gebruikelijke strategie
zou zijn om te proberen $!A \lor \lnot !A$ met $\Intro{\lor}$ te
!!{derive}. Daarvoor zouden we echter zonder aannamen hetzij $!A$,
hetzij $\lnot !A$ moeten !!{derive}, en dat is onmogelijk. \FalseCl{}
brengt uitkomst.
\begin{prooftree}
  \AxiomC{$\Discharge{\lnot(!A \lor \lnot !A)}{1}$}
  \DeduceC{$\lfalse$}
  \DischargeRule{\FalseCl}{1}
  \UnaryInfC{$!A \lor \lnot !A$}
\end{prooftree}
Nu zoeken we !!a{derivation} van $\lfalse$ uit $\lnot(!A \lor
\lnot !A)$. Omdat $\lfalse$ de conclusie van $\Elim{\lnot}$ is,
kunnen we die regel proberen:
\begin{prooftree}
  \AxiomC{$\Discharge{\lnot(!A \lor \lnot !A)}{1}$}
  \DeduceC{$\lnot !A$}
  \AxiomC{$\Discharge{\lnot(!A \lor \lnot !A)}{1}$}
  \DeduceC{$!A$}
  \RightLabel{\Elim{\lnot}}
  \BinaryInfC{$\lfalse$}
  \DischargeRule{\FalseCl}{1}
  \UnaryInfC{$!A \lor \lnot !A$}
\end{prooftree}
Onze strategie om !!a{derivation} van~$\lnot !A$ te vinden, vereist
een toepassing van~$\Intro{\lnot}$:
\begin{prooftree}
  \AxiomC{$\Discharge{\lnot(!A \lor \lnot !A)}{1}, \Discharge{!A}{2}$}
  \DeduceC{$\lfalse$}
  \DischargeRule{\Intro{\lnot}}{2}
  \UnaryInfC{$\lnot !A$}
  \AxiomC{$\Discharge{\lnot(!A \lor \lnot !A)}{1}$}
  \DeduceC{$!A$}
  \RightLabel{\Elim{\lnot}}
  \BinaryInfC{$\lfalse$}
  \DischargeRule{\FalseCl}{1}
  \UnaryInfC{$!A \lor \lnot !A$}
\end{prooftree}
Hier kunnen we $\lfalse$ gemakkelijk verkrijgen door
$\Elim{\lnot}$ toe te passen op de aanname
$\lnot(!A \lor \lnot !A)$ en op $!A \lor \lnot !A$. De laatste
formule volgt uit onze nieuwe aanname $!A$ met~$\Intro{\lor}$:
\begin{prooftree}
  \AxiomC{$\Discharge{\lnot(!A \lor \lnot !A)}{1}$}
  \AxiomC{$\Discharge{!A}{2}$}
  \RightLabel{\Intro{\lor}}
  \UnaryInfC{$!A \lor \lnot !A$}
  \RightLabel{\Elim{\lnot}}
  \BinaryInfC{$\lfalse$}
  \DischargeRule{\Intro{\lnot}}{2}
  \UnaryInfC{$\lnot !A$}
  \AxiomC{$\Discharge{\lnot(!A \lor \lnot !A)}{1}$}
  \DeduceC{$!A$}
  \RightLabel{\Elim{\lnot}}
  \BinaryInfC{$\lfalse$}
  \DischargeRule{\FalseCl}{1}
  \UnaryInfC{$!A \lor \lnot !A$}
\end{prooftree}
Rechts gebruiken we dezelfde strategie, behalve dat we $!A$ met
~\FalseCl verkrijgen:
\begin{prooftree}
  \AxiomC{$\Discharge{\lnot(!A \lor \lnot !A)}{1}$}
  \AxiomC{$\Discharge{!A}{2}$}
  \RightLabel{\Intro{\lor}}
  \UnaryInfC{$!A \lor \lnot !A$}
  \RightLabel{\Elim{\lnot}}
  \BinaryInfC{$\lfalse$}
  \DischargeRule{\Intro{\lnot}}{2}
  \UnaryInfC{$\lnot !A$}
  \AxiomC{$\Discharge{\lnot(!A \lor \lnot !A)}{1}$}
  \AxiomC{$\Discharge{\lnot !A}{3}$}
  \RightLabel{\Intro{\lor}}
  \UnaryInfC{$!A \lor \lnot !A$}
  \RightLabel{\Elim{\lnot}}
  \BinaryInfC{$\lfalse$}
  \DischargeRule{\FalseCl}{3}
  \UnaryInfC{$!A$}
  \RightLabel{\Elim{\lnot}}
  \BinaryInfC{$\lfalse$}
  \DischargeRule{\FalseCl}{1}
  \UnaryInfC{$!A \lor \lnot !A$}
\end{prooftree}
\end{ex}

\begin{prob}
Geef !!{derivation}s die het volgende aantonen:
\begin{enumerate}
\item $!A \land (!B \land !C) \Proves (!A \land !B) \land !C$.
\item $!A \lor (!B \lor !C) \Proves (!A \lor !B) \lor !C$.
\item $!A \lif (!B \lif !C) \Proves !B \lif (!A \lif !C)$.
\item $!A \Proves \lnot\lnot !A$.
\end{enumerate}
\end{prob}

\begin{prob}
Geef !!{derivation}s die het volgende aantonen:
\begin{enumerate}
\item $(!A \lor !B) \lif !C \Proves !A \lif !C$.
\item $(!A \lif !C) \land (!B \lif !C) \Proves (!A \lor !B) \lif !C$.
\item $\Proves \lnot(!A \land \lnot !A)$.
\item $!B \lif !A \Proves \lnot !A \lif \lnot !B$.
\item $\Proves (!A \lif \lnot !A) \lif \lnot !A$.
\item $\Proves \lnot(!A \lif !B) \lif \lnot !B$.
\item $!A \lif !C \Proves \lnot (!A \land \lnot !C)$.
\item $!A \land \lnot !C \Proves \lnot (!A \lif !C)$.
\item $!A \lor !B, \lnot !B \Proves !A$.
\item $\lnot !A \lor \lnot !B \Proves \lnot(!A \land !B)$.
\item $\Proves (\lnot !A \land \lnot !B) \lif\lnot(!A \lor !B)$.
\item $\Proves \lnot(!A \lor !B) \lif (\lnot !A \land \lnot !B)$.
\end{enumerate}
\end{prob}

\begin{prob}
Geef !!{derivation}s die het volgende aantonen:
\begin{enumerate}
\item $\lnot(!A \lif !B) \Proves !A$.
\item $\lnot(!A \land !B) \Proves \lnot !A \lor \lnot !B$.
\item $!A \lif !B \Proves \lnot !A \lor !B$.
\item $\Proves \lnot \lnot !A \lif !A$.
\item $!A \lif !B, \lnot !A \lif !B \Proves !B$.
\item $(!A \land !B) \lif !C \Proves (!A \lif !C) \lor (!B \lif !C)$.
\item $(!A \lif !B) \lif !A \Proves !A$.
\item $\Proves (!A \lif !B) \lor (!B \lif !C)$.
\end{enumerate}
(Hiervoor is steeds de regel $\FalseCl$ nodig.)
\end{prob}
\end{document}
